
%(BEGIN_QUESTION)
% Copyright 2006, Tony R. Kuphaldt, released under the Creative Commons Attribution License (v 1.0)
% This means you may do almost anything with this work of mine, so long as you give me proper credit

Beskriv hvordan en {\it posisjoner} kan brukes til å endre den iboende karakteristikken til en reguleringsventil, som et alternativ til å bytte ventiltrim.

\vskip 10pt

Forklar også hvorfor bruk av en posisjoner til dette kanskje ikke gir best resultat.

\vskip 20pt \vbox{\hrule \hbox{\strut \vrule{} {\bf Forslag til sokratisk diskusjon} \vrule} \hrule}

\begin{itemize}
\item{} Beskriv et realistisk scenario der noen ønsker å endre karakteristikken til en reguleringsventil, og hvordan de kan vite at dette er riktig tiltak.
\end{itemize}

\underbar{file i01385}
%(END_QUESTION)





%(BEGIN_ANSWER)

Posisjonerere kan med vilje gjøres ikke-lineære, slik at ventilspindelens (eller akselens) posisjon ikke er lineært relatert til reguleringssignalet slik den normalt ville vært. Ved bruk av spesialformede kammer og andre mekaniske komponenter i posisjonerens tilbakekoblingsmekanisme kan man oppnå nesten hvilken som helst åpningskarakteristikk fra én og samme trimtype.

\vskip 10pt

Dette betyr likevel ikke at man kan få {\it nøyaktig} equal-percentage-åpningskarakteristikk fra en ventiltrim som i utgangspunktet er quick-opening. Fordi posisjoneren må bevege en quick-opening-mekanisme med ekstrem presisjon i nedre del av arbeidsområdet for å etterligne equal percentage (altså tvinge en ventil som åpner aggressivt til å åpne mer forsiktig når reguleringssignalet øker), kan denne metoden ikke garantere god nøyaktighet over hele ventilens arbeidsområde. Den har likevel en klar fordel: man kan endre ventilens åpningskarakteristikk uten å demontere ventilen eller ta den ut av røranlegget!

%(END_ANSWER)





%(BEGIN_NOTES)


%INDEX% Final Control Elements, valve: characterization

%(END_NOTES)

